\corrigesChapitre{7}{Combustions et conservation de la masse}
\par\noindent\begin{minipage}{\linewidth}
\begin{corrige}{1}{Deux rôles}
\begin{questions}
\item Le charbon joue le rôle de combustible ; le dioxygène, de comburant.
\item Non : l’air est un mélange ; on distingue ses constituants et le dioxygène consommé.
\end{questions}
\end{corrige}
\end{minipage}
\par\noindent\begin{minipage}{\linewidth}
\begin{corrige}{2}{Les produits du méthane}
\begin{questions}
\item Dioxyde de carbone et eau.
\item Oui, une combustion peut transférer de l’énergie notamment thermiquement et par rayonnement.
\end{questions}
\end{corrige}
\end{minipage}
\par\noindent\begin{minipage}{\linewidth}
\begin{corrige}{3}{Une somme de masses}
\begin{questions}
\item $7\unite g$.
\item La mesure doit inclure toute la matière et aucun produit ne doit quitter le système pesé.
\end{questions}
\end{corrige}
\end{minipage}
\par\noindent\begin{minipage}{\linewidth}
\begin{corrige}{4}{Deux gaz à distinguer}
\begin{questions}
\item CO comporte un O pour un C ; CO2 en comporte deux.
\item Ni l’un ni l’autre ne doit être identifié en le sentant ; le CO est notamment inodore et très toxique.
\end{questions}
\end{corrige}
\end{minipage}
\par\noindent\begin{minipage}{\linewidth}
\begin{corrige}{5}{Une condition manquante}
\begin{questions}
\item Combustible, comburant et énergie d’activation.
\item Uniquement dans une activité encadrée par le professeur, avec le dispositif et les protections prévus.
\end{questions}
\end{corrige}
\end{minipage}
\par\noindent\begin{minipage}{\linewidth}
\begin{corrige}{6}{Compter les atomes}
\begin{questions}
\item Quatre H et quatre O de chaque côté.
\item Elles contiennent les quatre H du méthane et deux des quatre O ; les deux autres O sont dans le dioxyde de carbone.
\end{questions}
\pcFigureEx{figures/cor-reaction-methane.pdf}{Un C, quatre H et quatre O se retrouvent de chaque côté ; le nombre de molécules n’est pas imposé constant.}
\end{corrige}
\end{minipage}
\par\noindent\begin{minipage}{\linewidth}
\begin{corrige}{7}{Une masse manquante}
\begin{questions}
\item $22-6=16\unite g$.
\item Tous les produits sont pris en compte et il n’y a ni perte ni autre échange de matière dans le bilan.
\end{questions}
\par\smallskip\noindent\textbf{Pour comprendre.}\enspace Conservation de masse dans le bilan complet : $m_1+m_2=m_{\text{produit}}$. D’où $m_2=22-6=16\unite g$. Le calcul suppose les deux réactifs entièrement consommés et aucun échange non compté.\par
\end{corrige}
\end{minipage}
\par\noindent\begin{minipage}{\linewidth}
\begin{corrige}{8}{Le produit est plus lourd}
\begin{questions}
\item $44-12=32\unite g$.
\item Non : les $32\unite g$ supplémentaires viennent du dioxygène, qui doit être inclus dans les réactifs.
\end{questions}
\par\smallskip\noindent\textbf{Pour comprendre.}\enspace $m_{\text{dioxygène}}=44-12=32\unite g$. Le produit est plus lourd que le carbone seul parce qu’il contient aussi la matière apportée par le dioxygène.\par
\end{corrige}
\end{minipage}
\par\noindent\begin{minipage}{\linewidth}
\begin{corrige}{9}{Un montage ouvert}
\begin{questions}
\item Non : le système limité à la bougie échange de la matière avec l’air.
\item Le dioxygène consommé et les produits gazeux qui s’échappent, en plus de la cire et des éventuels résidus.
\end{questions}
\pcFigureEx{figures/cor-bougie-ouverte.pdf}{Le bilan doit inclure le dioxygène entrant et les produits gazeux sortants.}
\end{corrige}
\end{minipage}
\par\noindent\begin{minipage}{\linewidth}
\begin{corrige}{10}{Un bilan proportionnel}
\begin{questions}
\item On multiplie par 3 : $24\unite g$ de dioxygène et $33\unite g$ de produit.
\item Non : il concerne le dioxygène pur consommé ; l’air comporte aussi d’autres gaz.
\end{questions}
\par\smallskip\noindent\textbf{Pour comprendre.}\enspace Passer de 3 g à 9 g de carbone multiplie les masses du bilan par trois : 24 g de dioxygène et 33 g de produit. Le bilan chimique utilise O2 ; l’air contient aussi d’autres gaz.\par
\end{corrige}
\end{minipage}
\par\noindent\begin{minipage}{\linewidth}
\begin{corrige}{11}{Une affirmation à corriger}
\begin{questions}
\item Elle peut produire du monoxyde de carbone, de la suie et d’autres produits selon les conditions.
\item Le monoxyde de carbone est inodore ; son absence d’odeur n’atteste pas l’absence de danger.
\end{questions}
\end{corrige}
\end{minipage}
\par\noindent\begin{minipage}{\linewidth}
\begin{corrige}{12}{Une mesure presque égale}
\begin{questions}
\item Non : un écart de $0{,}1\unite g$ doit être discuté avec la précision et la variabilité des mesures.
\item La stabilité de la balance, l’étanchéité, les conditions de pesée et la reproductibilité.
\end{questions}
\end{corrige}
\end{minipage}
\par\noindent\begin{minipage}{\linewidth}
\begin{corrige}{13}{Un réactif en excès}
\begin{questions}
\item Les $6\unite g$ de C demandent $16\unite g$ de dioxygène ; il reste $8\unite g$ de dioxygène.
\item $6+16=22\unite g$ de produit ; $22+8=30\unite g$ est la masse finale totale.
\end{questions}
\par\smallskip\noindent\textbf{Pour comprendre.}\enspace Les 6 g de carbone demandent $6\times8/3=16\unite g$ de dioxygène. Sur 24 g, il en reste 8 g ; le produit formé pèse $6+16=22\unite g$. Contrôle : $22+8=6+24=30\unite g$.\par
\pcFigureEx{figures/cor-bilan-carbone-6.pdf}{Séparer matière transformée et réactif restant ; la masse finale totale comprend les deux.}
\end{corrige}
\end{minipage}
\par\noindent\begin{minipage}{\linewidth}
\begin{corrige}{14}{Émissions comparées}
\begin{questions}
\item A : $1\,000\unite g$ ; B : $400\unite g$.
\item Non : cette comparaison porte sur ce seul indicateur et le périmètre de la donnée, sans inclure nécessairement fabrication, ressources et fin de vie.
\end{questions}
\end{corrige}
\end{minipage}
\par\noindent\begin{minipage}{\linewidth}
\begin{corrige}{15}{Conserver quoi ?}
\begin{questions}
\item Non : $2\,\ce{H2}+\ce{O2}\rightarrow2\,\ce{H2O}$ passe de trois molécules à deux dans le bilan.
\item Le nombre d’atomes de chaque élément et la masse totale du système matériel complet.
\end{questions}
\end{corrige}
\end{minipage}
\par\noindent\begin{minipage}{\linewidth}
\phantomsection\label{cor-defi-7}\pcSolutionDefi{7}{Le bilan sous contrôle}\begin{questions}
\item Les $9\unite g$ de C demanderaient $24\unite g$ de dioxygène : le dioxygène disponible est insuffisant.
\item Les $20\unite g$ correspondent à $20\times3/8=7{,}5\unite g$ de C. Il reste $1{,}5\unite g$ ; on forme $27{,}5\unite g$ de produit dans ce modèle.
\item $27{,}5+1{,}5=29\unite g$, comme $9+20$. Le calcul utilise un bilan de combustion complète idéal ; des conditions réelles pauvres en dioxygène peuvent produire d’autres espèces.
\end{questions}
\pcFigureEx{figures/cor-bilan-carbone-9.pdf}{Le dioxygène disponible limite ici le modèle de transformation complète ; inclure le carbone restant dans la masse finale.}
\end{minipage}
\par\noindent\begin{minipage}{\linewidth}
\phantomsection\label{cor-activite-7}\pcSousTitre{Pistes pour l’investigation}
Dans le dispositif fermé adapté, la masse totale avant et après est compatible avec une conservation à la précision de mesure. Dans un système ouvert, les entrées et sorties gazeuses expliquent une variation de la partie pesée. Le schéma du système doit donc précéder le calcul. On conserve les nombres d’atomes de chaque élément, pas nécessairement le nombre de molécules.\par\end{minipage}
